\section*{Resolução da Questão 5: Teoria dos Jogos (Minimax e Alpha-Beta)}

\subsection*{A) Fundamentos Teóricos Adversariais}

\begin{enumerate}[label=\roman*)]
    \item \textbf{Meta do Minimax:} Consiste em extrair o roteiro inquebrável para um protagonista, assegurando o ganho mínimo em face de um nêmese perfeito que quer arruiná-lo.
    \item \textbf{Dicotomia MAX/MIN:} Vértices MAX mapeiam o turno do nosso agente, filtrando o filho com pontuação mais gorda. Vértices MIN simulam o algoz adversário, punindo-nos ao escolher a rota com a pontuação mais miserável possível.
    \item \textbf{Carga de Utilidade:} A tradução quantitativa da vitória, derrota ou saldo do jogo contida estritamente nas folhas da árvore final.
    \item \textbf{Fluxo Retroativo:} Como a partida é lida do futuro para o presente, o sistema despenca até os nós terminais e canaliza as notas de volta para a cúpula, onde as hierarquias MAX/MIN engolem as piores ou melhores opções.
    \item \textbf{A Premissa da Perfeição:} É o postulado de que o Inimigo possui onisciência computacional e jamás sofrerá deslizes ou blefes falhos; ele jogará para nos liquidar.
    \item \textbf{Cirurgia Alpha-Beta:} Uma tática extrema de corte que poda galhos redundantes cujo acesso matemático já se prova impossível sem sequer lê-los, poupando uma massa colossal de memória e processamento.
\end{enumerate}

\subsection*{B) Avaliação do Eixo Minimax}

\textbf{Composição Analítica Final:} Após escavar todo o tronco C-D-E-H, a árvore destila a verdade matemática do jogo.

\begin{table}[H]
\centering
\begin{tabular}{|c|c|c|}
\hline
\textbf{Vértice} & \textbf{Polo} & \textbf{Avaliação Consolidada} \\ \hline
A & MAX & 6 \\ \hline
B & MIN & 5 \\ \hline
C & MAX & 5 \\ \hline
D & MAX & 9 \\ \hline
E & MIN & 0 \\ \hline
F & MAX & 2 \\ \hline
G & MAX & 0 \\ \hline
H & MIN & 6 \\ \hline
I & MAX & 7 \\ \hline
J & MAX & 6 \\ \hline
\end{tabular}
\end{table}

\begin{itemize}
    \item \textbf{Decisão de MAX na Raiz:} Enveredou para a ala direita, consolidando a utilidade global de \textbf{6}.
    \item \textbf{Trajeto Consagrado:} $A \rightarrow H \rightarrow J \rightarrow Folha(6)$.
\end{itemize}

\textbf{Modelagem da Árvore de Decisão:}
\textit{Caminho rei assinalado em amarelo forte (\textcolor{yellow}{Cor}).}

\begin{center}
\begin{forest}
  for tree={
    circle,
    draw,
    thick,
    minimum size=8mm,
    inner sep=1pt,
    s sep=5mm,
    l sep=10mm,
    font=\small
  }
  [$A(6)$, fill=yellow!50
    [$B(5)$
      [$C(5)$ [3] [5]]
      [$D(9)$ [6] [9]]
    ]
    [$E(0)$
      [$F(2)$ [1] [2]]
      [$G(0)$ [0] [-1]]
    ]
    [$H(6)$, fill=yellow!50
      [$I(7)$ [7] [4]]
      [$J(6)$, fill=yellow!50
        [5] 
        [6, fill=yellow!50]
      ]
    ]
  ]
\end{forest}
\end{center}

\subsection*{C) Execução da Poda Alpha-Beta}

\begin{table}[H]
\centering
\begin{tabular}{|c|c|c|c|c|}
\hline
\textbf{Ciclo} & \textbf{Avaliado} & \textbf{Apoio $\alpha$} & \textbf{Apoio $\beta$} & \textbf{Corte Executado?} \\ \hline
1 & C & 3 & $\infty$ & Não \\ \hline
2 & C & 5 & $\infty$ & Não \\ \hline
3 & B & $-\infty$ & 5 & Não \\ \hline
4 & D & 6 & 5 & \textbf{Sim} ($\alpha \ge \beta$) \\ \hline
5 & B & $-\infty$ & 5 & Não \\ \hline
6 & A & 5 & $\infty$ & Não \\ \hline
7 & F & 5 & $\infty$ & Não \\ \hline
8 & F & 5 & $\infty$ & Não \\ \hline
9 & E & 5 & 2 & \textbf{Sim} ($\alpha \ge \beta$) \\ \hline
10 & A & 5 & $\infty$ & Não \\ \hline
11 & I & 7 & $\infty$ & Não \\ \hline
12 & I & 7 & $\infty$ & Não \\ \hline
13 & H & 5 & 7 & Não \\ \hline
14 & J & 5 & 7 & Não \\ \hline
15 & J & 6 & 7 & Não \\ \hline
16 & H & 5 & 6 & Não \\ \hline
17 & A & 6 & $\infty$ & Não \\ \hline
\end{tabular}
\end{table}

A dinâmica fulminou ramos inteiros:
\begin{itemize}
    \item Em **D**, no instante que a folha 6 inflou o $\alpha$ local de D, o $\beta$ de B já estava sedimentado em 5. Ninguém cederia, e o nó virou lixo lógico.
    \item Em **E**, a mesma discrepância dizimou inteiramente o acesso a **G**, preservando o fluxo de dados em um nível de processamento de apenas 18 nós (contrastando com os 22 absolutos da busca completa).
\end{itemize}

\subsection*{D) Experimento: Distorção de Valores nas Folhas}

Refazendo o roteiro, injetamos instabilidade:
\begin{itemize}
    \item Folha 1 de F passou de 1 para \textbf{12}.
    \item Folha 2 de F passou de 2 para \textbf{18}.
    \item Folha 1 de I despencou de 7 para \textbf{-10}.
\end{itemize}

\textbf{Rastreio Dinâmico (Alpha-Beta Pós-Mutação):}

\begin{table}[H]
\centering
\begin{tabular}{|c|c|c|c|c|}
\hline
\textbf{Ciclo} & \textbf{Vértice} & \textbf{$\alpha$} & \textbf{$\beta$} & \textbf{Podou?} \\ \hline
... & ... & ... & ... & ... \\ \hline
6 & A & 5 & $\infty$ & Não \\ \hline
7 & F & 12 & $\infty$ & Não \\ \hline
8 & F & 18 & $\infty$ & Não \\ \hline
9 & E & 5 & 18 & Não \\ \hline
10 & G & 5 & $\infty$ & Não \\ \hline
11 & G & 5 & $\infty$ & Não \\ \hline
12 & E & 5 & 0 & \textbf{Sim} ($\alpha \ge \beta$) \\ \hline
13 & A & 5 & $\infty$ & Não \\ \hline
14 & I & -10 & $\infty$ & Não \\ \hline
15 & H & 5 & -10 & \textbf{Sim} ($\alpha \ge \beta$) \\ \hline
16 & A & 5 & $\infty$ & Não \\ \hline
\end{tabular}
\end{table}

\textbf{Conclusões Experimento Efeito Cascata:}
A presença venenosa da folha "-10" em I contaminou mortalmente o antigo caminho ouro de H. O adversário (MIN) em H viu o -10 e travou esse teto catastrófico de perdas para nós. Com $\beta = -10$, ocorreu uma poda massiva em H contra a segurança do $\alpha = 5$ da raiz A. Descartado o eixo H, fomos jogados à mercê da segunda melhor alternativa, alterando o desfecho mundial da partida inteira para o lucro mediano de \textbf{5} originado pelo bloco B. A heurística limitada não foi capaz de antecipar um desastre e a matemática fria salvou o jogador de perder com lucro negativo.
